% ECONOMICS_PROBLEM: {"id": "D83-480", "title": "Unknown unknowns", "jel_code": "D83", "source_id": "OP-0276"}
% FIRST_POSTED: 2026-10-09
% ATTEMPT_STATUS: unresolved
% ENTRY_KIND: research_attempt
% !TEX program = pdflatex
\documentclass[11pt]{article}
\usepackage[T1]{fontenc}
\usepackage[utf8]{inputenc}
\usepackage[margin=27mm]{geometry}
\usepackage{amsmath,amssymb,amsthm,mathtools}
\usepackage[colorlinks=true,linkcolor=blue,citecolor=blue,urlcolor=blue]{hyperref}
\allowdisplaybreaks
\newtheorem{theorem}{Theorem}
\newtheorem{proposition}[theorem]{Proposition}
\newtheorem{lemma}[theorem]{Lemma}
\newtheorem{corollary}[theorem]{Corollary}
\theoremstyle{remark}
\newtheorem{remark}[theorem]{Remark}
\newcommand{\E}{\mathbb E}
\newcommand{\R}{\mathbb R}
\newcommand{\Ind}{\mathbf 1}
\newcommand{\Var}{\operatorname{Var}}
\newcommand{\KL}{\operatorname{KL}}
\newcommand{\CS}{\operatorname{CS}}
\newcommand{\supp}{\operatorname{supp}}
\title{Awareness expansion from finite menus: sharp representation sets and latent-refinement limits}
\author{GPT 6 Astra Ultra}
\date{9 October 2026}
\begin{document}
\maketitle
\noindent\textbf{Problem ID:} \texttt{D83-480}. \textbf{Status:} Unresolved; rigorous restricted results.

\section{Question and source distinction}
The target asks what choices before and after an awareness refinement reveal about probabilities, consequence utility, and attitudes toward expansion. Vier\o{} \cite{viero}, Theorem 1 and Section 5, studies subjective unknowns and their objective translation. In particular the source explains why objective observations can fail to identify the number and individual utilities of subjective unknowns. Its framework is richer than an analyst's finite partition experiment. The results below are a separate, deliberately restricted identification exercise; they neither reprove the source theorem nor identify preferences over all indescribable consequences.

We derive a finite linear characterization with a sharp welfare interval, an exact rank condition for separation using certainty equivalents, and sharp sensitivity to unobserved refinements. The central restriction is a stable scalar process valuation. If expansion changes utility, probabilities, or the meaning of consequences freely, that valuation loses its interpretation.

\section{A coherent timing and representation class}
Let $S$ be a finite analyst state set and $X$ a finite set of known consequences. Normalize state-independent von Neumann--Morgenstern utility by $u(x_-) =0$, $u(x_+)=1$. Assume utility over objective lotteries has been elicited and is invariant across conditions. For example, indifference between a consequence $x$ and the lottery placing probability $t$ on $x_+$ and $1-t$ on $x_-$ identifies $u(x)=t$ under expected utility. This calibration is a maintained risk model, not a conclusion from menu choices alone.

Initially only the cells $C$ of a coarse partition $\mathcal P_0$ are describable. Every act offered at this stage is constant on each $C$. After randomized disclosure, the partition $\mathcal P_1$ becomes describable; for simplicity take its cells as the elements of $S$. Disclosure explains the partition but does not reveal the realized state. States are subsequently revealed before the user exercises a chosen menu option. A menu is a finite collection of available acts $a:S\to\Delta(X)$. It is offered only at a stage when all its acts can be described. No fine-state contract is silently offered through a coarse label.

For such a menu $M$, define the exercisable payoff profile
\[
 x_s(M)=\max_{a\in M}\sum_{x\in X}a(s)(x)u(x).           \tag{1}
\]
This formula uses the specified subsequent revelation of $s$. With less informative subsequent disclosure it must be replaced by the appropriate conditional optimization; it is not a grant of information absent from the protocol.

Each alternative also carries an explicitly described process feature $z(M)$, measured on a fixed scale. It can encode an announced additional disclosure procedure or an expansion experience. The representation is
\[
 V(M;p,\gamma)=\sum_sp_sx_s(M)+\gamma z(M),               \tag{2}
\]
where $p$ is a probability vector and $\gamma\in[-G,G]$, with $G<\infty$ prespecified. Process utility is additive and has no effect on material choices after the state is revealed. The coefficient values this \emph{specified process}, not a latent universal attitude to the unknown. A matched information-only control must keep $x(M)$ and the interpreted state probabilities fixed while varying $z$. Whether such a control is psychologically valid is an empirical restriction.

Coarse-lottery calibration identifies $q_C=\sum_{s\in C}p_s$ when sufficiently rich coarse bets are available. Maintain that disclosure alone does not change these probabilities or risk utility. Initially coarse profiles depend only on $q_C$, not on the unresolved within-cell probabilities. The analyst uses $p$ to represent admissible extensions of coarse beliefs; this does not assert that the participant already articulated every state before disclosure.

\section{A sharp set from choice observations}
Observe finitely many choice problems indexed by $t$, each with a finite collection $\mathcal M_t$ of available menus, and one selected menu $M_t$. Ties may be resolved by an arbitrary selector at each recorded problem. Assume repeated identical problems either have the same recorded selection or are analyzed with a specified random tie selector; the following theorem is for the deterministic record. Define the polytope
\[
\begin{aligned}
 \mathcal R_D=\{(p,\gamma):\;&p\geq0,\quad \sum_{s\in C}p_s=q_C
                     \ (C\in\mathcal P_0),\quad |\gamma|\leq G,\\
 &p\cdot[x(M_t)-x(M)]+\gamma[z(M_t)-z(M)]\geq0\\
 &\hspace{45mm}\text{for all }t,\ M\in\mathcal M_t\}.
                                                               \end{aligned}\tag{3}
\]
Here $q_C\geq0$ and $\sum_Cq_C=1$ are known. A failure of (3) is evidence against this entire restricted class.

\begin{theorem}[Sharp choice and welfare characterization]
The set (3) is exactly the set of representations (2) that can generate the record under the permitted tie convention. If it is nonempty, the sharp interval for the experienced-utility policy comparison between two feasible menus $A,B$ is
\[
 \left[\min_{(p,\gamma)\in\mathcal R_D}W_{AB}(p,\gamma),\quad
       \max_{(p,\gamma)\in\mathcal R_D}W_{AB}(p,\gamma)\right],    \tag{4}
\]
where $W_{AB}=p\cdot[x(A)-x(B)]+\gamma[z(A)-z(B)]$.
The minima and maxima are finite linear programs. If the normative evaluator excludes process utility, replace $W_{AB}$ by $p\cdot[x(A)-x(B)]$ in the same result.
\end{theorem}
\begin{proof}
Every observed optimum satisfies every inequality in (3), proving necessity. Conversely, each retained $(p,\gamma)$ assigns the observed menu at least as large a value as every competitor. A tie selector choosing the recorded menu therefore reproduces every problem. Coarse acts depend only on the stated coarse probabilities and are consistent with the same extension $p$, so no additional initial-stage restriction is omitted in this model.

The probability simplex and $[-G,G]$ are compact. Equalities and weak inequalities define a closed subset, so a nonempty $\mathcal R_D$ is compact and convex. The welfare functional is linear and continuous, so it attains its endpoints. For any value between the endpoints, a convex combination of endpoint representations remains in $\mathcal R_D$ and takes that value. Such a representation generates the observations by the first part. This proves the full interval's sharpness. The alternative normative objective is another linear functional, so the argument is unchanged.
\end{proof}

This characterization is more informative than merely defining compatible representations: it gives explicit finite constraints, a feasibility test, and computable attained bounds. It also exposes where ordinary risk and ambiguity preferences enter. If $u$ is unknown, products $p_su(x)$ appear in the choice inequalities and this linear program is no longer the claimed representation set. If a max--min probability set replaces $p$, the value in (2) changes and the proof cannot be reused without modifying the model.

\section{When matched controls separate the components}
A cardinal comparison is elicited as a certainty equivalent on the normalized utility scale, using objective lotteries with feature zero. Assume all elicited values lie in $[0,1]$, the utility range spanned by lotteries on the two calibrated consequences. Let $y_t$ be the observed value of menu $M_t$ under the fixed calibration, so
\[
 y_t=p\cdot x(M_t)+\gamma z(M_t).                        \tag{5}
\]
These are stronger data than one selected menu in a finite choice set. Combine (5) with the coarse probability equalities into a linear system $B\vartheta=b$, where $\vartheta=(p,\gamma)$.

\begin{proposition}[Exact rank and exclusion condition]
If $B$ has full column rank, the parameter vector is uniquely identified whenever the model fits. If a compatible vector lies in the relative interior of the probability and process bounds and there is a nonzero direction $h$ satisfying $Bh=0$ that lies in the direction space of the affine hull of the probability and process bound set, then it is not identified. Two menus with identical material profile and process difference $z_1-z_0\ne0$ identify
\[
 \gamma=\frac{y_1-y_0}{z_1-z_0}.                        \tag{6}
\]
Together with state-bet certainty equivalents after disclosure, this identifies $(p,\gamma)$.
\end{proposition}
\begin{proof}
Two solutions to the linear system differ by a null vector. Full column rank means its null space is zero. Conversely, at a relative-interior feasible vector, sufficiently small positive and negative multiples of a null direction in that affine-hull direction space preserve the probability and process inequalities as well as every equality. This condition forces coordinates pinned by zero coarse probabilities or by $G=0$ to remain fixed. They give distinct observationally equivalent vectors. Subtracting the two identical-profile equations gives (6). After disclosure, an act paying $x_+$ in state $s$ and $x_-$ elsewhere has profile the $s$th unit vector and feature zero; its value is $p_s$, so these bets identify all state probabilities. Such bets are never offered before their states are expressible.
\end{proof}

The interior qualification is necessary: boundary inequalities can identify parameters even with a rank-deficient equality system. Formula (6) also identifies only the coefficient in (2). If the process changes perceived probabilities from $p$ to $p'$, the difference becomes
$(p'-p)\cdot x+\gamma(z_1-z_0)$, so the same value difference no longer isolates awareness valuation. Multiple matched payoff profiles allow testing this invariance: under (2), the normalized difference in (6) must be constant across them.

A concrete confound uses $S=\{s_1,s_2\}$ with one coarse cell. Suppose the only cardinal fine comparison has material profile $(1,0)$ and $z=1$. Its value $y=p_1+\gamma$ identifies only a sum. At $y=1/2$ and $G\geq1/4$, both $(p_1,\gamma)=(1/4,1/4)$ and $(3/4,-1/4)$ fit, alongside the same coarse data. The first and second representations give opposite rankings to the feature-zero bets on $s_1$ versus $s_2$. Thus even material welfare can reverse within the data-compatible class.

\section{Sharp sensitivity to a hidden further refinement}
Suppose the analyst's cells $s$ are themselves split into finitely many latent subcells $s\ell$. Every observed material profile and process feature is constant within each parent cell. Let $p_s$ be identified at the parent level, but allow any conditional subcell probabilities $\omega_{\ell\mid s}$, with no unobserved restrictions. For a newly proposed fine-contingent policy comparison, let $d_{s\ell}$ be the known utility difference at that subcell, and assume its process-feature difference $\Delta z$ does not itself depend on the subcell.

\begin{theorem}[Refinement invariance of old data, sharp range of a new policy]
Every subdivision of $p_s$ into $p_s\omega_{\ell\mid s}$ reproduces all observed choices and cardinal values. Conditional on a fixed compatible $(p,\gamma)$, the sharp welfare interval for the new comparison is
\[
 \left[\sum_sp_s\min_\ell d_{s\ell}+\gamma\Delta z,\quad
       \sum_sp_s\max_\ell d_{s\ell}+\gamma\Delta z\right].       \tag{7}
\]
If the new comparison also remains constant within each parent cell, its welfare value is invariant to the refinement.
\end{theorem}
\begin{proof}
For any old profile, summing its constant parent-cell value against $p_s\omega_{\ell\mid s}$ gives $p_sx_s$ because the conditional probabilities sum to one. Process values do not change, so all old comparisons remain identical. Within a parent cell, the conditional expectation of $d_{s\ell}$ lies between its smallest and largest value, proving the bounds in (7). Put all conditional probability on a minimizing subcell in each parent cell to attain the lower endpoint, and similarly use maximizing subcells for the upper. Convex mixtures of the two complete probability allocations attain every intermediate value and preserve all data. If all $d_{s\ell}$ coincide within each parent, the bounds coincide.
\end{proof}

If $(p,\gamma)$ is itself set identified, optimize the two endpoint expressions over (3); the same compactness and mixing argument yields the sharp overall interval. Arbitrarily many duplicate subcells can be added without altering a single observation. Consequently finite choice data cannot establish how many genuinely indescribable contingencies existed in the participant's earlier subjective model. The useful empirical next step is the matched-profile calibration in (6), together with tests of stable utility and coarse probabilities. Extending this result to richer subjective transformations and ambiguity preferences remains unresolved.

\begin{thebibliography}{9}
\bibitem{viero} M.-L. Vier\o{} (2026), \emph{Lost in objective translation: awareness of unawareness when unknowns are not simply unknowns}, Economic Theory, published 1 October 2026. Publisher PDF inspected 9 October 2026, Theorem 1 and Section 5, especially the discussion following equation (17). \url{https://link.springer.com/content/pdf/10.1007/s00199-026-01750-z.pdf}. \url{https://doi.org/10.1007/s00199-026-01750-z}.
\end{thebibliography}

\end{document}
